\documentclass[11pt,a4paper]{article}

% ---------- packages ----------
\usepackage[utf8]{inputenc}
\usepackage[T1]{fontenc}
\usepackage{lmodern}
\usepackage{amsmath,amssymb,amsthm}
\usepackage{mathtools}
\usepackage{physics}
\usepackage{bm}
\usepackage{geometry}
\geometry{margin=2.5cm}
\usepackage{enumitem}
\usepackage{booktabs}
\usepackage{caption}
\usepackage[dvipsnames]{xcolor}
\usepackage{hyperref}
\hypersetup{colorlinks,linkcolor=MidnightBlue,citecolor=OliveGreen,urlcolor=NavyBlue}
\usepackage{cleveref}
\usepackage{tikz}
\usetikzlibrary{calc,decorations.pathreplacing,patterns,arrows.meta}

% ---------- theorem styles ----------
\newtheorem{theorem}{Theorem}
\newtheorem{lemma}[theorem]{Lemma}
\newtheorem{proposition}[theorem]{Proposition}
\newtheorem{corollary}[theorem]{Corollary}
\theoremstyle{definition}
\newtheorem{definition}[theorem]{Definition}
\theoremstyle{remark}
\newtheorem{remark}[theorem]{Remark}
\newtheorem{claim}[theorem]{Claim}

% ---------- shortcuts ----------
\newcommand{\Sinc}{S_{\mathrm{inc}}}
\newcommand{\APD}{S_{\mathrm{ab}}}
\newcommand{\Tlay}{T_{\mathrm{lay}}}
\newcommand{\Tzero}{T_0}
\newcommand{\nhat}{\hat{\vb{n}}}
\newcommand{\khat}{\hat{\vb{k}}}
\newcommand{\ntilde}{\tilde{n}}
\newcommand{\SAR}{\ensuremath{\mathrm{SAR}}}
\newcommand{\psSAR}{\ensuremath{\mathrm{psSAR}_{10\mathrm{g}}}}
\newcommand{\rhom}{\rho_{\mathrm{m}}}

\title{%
  APD dominates $\psSAR$:\\
  a closed-form bound from energy conservation%
}
\author{Robin Wydaeghe (with Claude Opus 4.7)}
\date{2026-05-08}

\begin{document}
\maketitle

\begin{abstract}
This note proves a one-line inequality between the two ICNIRP basic
restrictions on local exposure: $\psSAR \le \APD/(\rhom L)$ in the planar
three-layer tissue model, where $\APD$ is the absorbed power density at
the body surface, $\rhom$ is tissue mass density, and $L = 21.5$~mm is
the IEC/IEEE 62704-1 cube side. At ICNIRP basic restrictions
($\APD \le 10$~W/m$^2$), this gives $\psSAR \le 0.47$~W/kg, below
the head and trunk basic restriction (2~W/kg) by a factor of $4.3$ and
the limb restriction (4~W/kg) by a factor of $8.6$. The bound is
energy-conservation, holds across $1$--$100$~GHz, and certifies $\psSAR$
compliance from any $\APD$ calculation. The closed-form chain that
delivers $\APD$ via the layered transmission $\Tlay$ therefore handles
both basic restrictions through a single surface integral.
\end{abstract}

\section{Setup}\label{sec:setup}

ICNIRP 2020 specifies two basic restrictions on local
exposure~\cite{ICNIRP2020}:

\begin{itemize}\setlength{\itemsep}{2pt}
  \item \textbf{Sub-6~GHz}: peak spatial-average SAR over a
        $10$~g cube of tissue, $\psSAR$. Limit $2$~W/kg head and trunk,
        $4$~W/kg limbs (general public, $6$-min average).
  \item \textbf{Above 6~GHz}: absorbed power density $\APD$ averaged
        over a $4$~cm$^2$ patch ($1$~cm$^2$ above $30$~GHz). Limit
        $10$~W/m$^2$ (general public).
\end{itemize}

\noindent
The cube quantity is operationalized by IEC/IEEE
62704-1~\cite{62704-1}, which prescribes axis-aligned cubes of side
$L = (m/\rhom)^{1/3} = 21.5$~mm at mass $m = 10$~g and
$\rhom = 1000$~kg/m$^3$, with the constraints (i) at least $90\%$ tissue
by mass and (ii) no face entirely in air. The compliance procedure
searches over valid placements and returns the maximum cube-averaged SAR
as $\psSAR$.

The two metrics measure related physics. $\APD$ is the time-averaged
real Poynting vector entering the body surface, integrated over a
patch. $\psSAR$ is the absorbed power per unit cube mass. Both depend on
the same underlying field; the difference is the spatial averaging
window (surface patch vs.\ volumetric cube) and the normalization
(per area vs.\ per mass).

The relationship below is energy-conservation. The cube can hold no
more absorbed power than enters it through its body-surface footprint.
Stating this carefully gives a closed-form bound that ties the two
basic restrictions together.

\section{The bound}\label{sec:bound}

\subsection{Result}

\begin{theorem}[APD dominates $\psSAR$]\label{thm:apd-dom}
For an axis-aligned $10$~g cube placed per IEC/IEEE 62704-1 on a planar
three-layer body, with cube z-axis aligned to the local outward body
normal $\nhat$, the peak spatial-average SAR satisfies
\begin{equation}\label{eq:apd-dom}
  \psSAR \;\le\; \frac{\APD}{\rhom\,L}\, ,
\end{equation}
where $\APD$ is the absorbed power density at the cube's body footprint.
For ICNIRP cube parameters, $\APD \le 10$~W/m$^2$ implies
$\psSAR \le 0.465$~W/kg.
\end{theorem}

\subsection{Proof}

The proof is two lines of energy conservation. By definition, $\APD$
is the time-averaged absorbed flux entering tissue per unit body-surface
area. Total absorbed power inside the cube cannot exceed total flux
entering through the cube's body-surface footprint of area $A$:
\begin{equation}\label{eq:energy-cons}
  P_{\mathrm{abs, cube}} \;\le\; \APD \cdot A\, .
\end{equation}
For a cube axis-aligned to $\nhat$, the body plane cuts perpendicular to
the cube z-axis and the intersection is the constant $L \times L$
cross-section: $A = L^2$. Cube mass is $m = \rhom L^3$. Dividing
\cref{eq:energy-cons} by $m$,
\begin{equation}
  \psSAR \;=\; \frac{P_{\mathrm{abs, cube}}}{m}
  \;\le\; \frac{\APD\,L^2}{\rhom L^3}
  \;=\; \frac{\APD}{\rhom L}\, ,
\end{equation}
which is~\cref{eq:apd-dom}. Equality requires every joule entering
through the footprint to deposit inside the cube, which holds in the
strong-attenuation limit $\alpha L \to \infty$ where $\alpha$ is the
field amplitude decay rate. \qed

\subsection{Numerical evaluation at ICNIRP limits}

With $\rhom = 1000$~kg/m$^3$ and $L = 21.5$~mm,
$\rhom L = 21.5$~kg/m$^2$. The bound becomes
\begin{equation}\label{eq:apd-dom-num}
  \psSAR \;\le\; \frac{\APD}{21.5~\mathrm{kg/m^2}}
  \;\approx\; 0.0465 \,\APD~\mathrm{W/kg}~~\text{per W/m}^2\, .
\end{equation}
At the ICNIRP basic restriction $\APD = 10$~W/m$^2$, this gives
$\psSAR \le 0.465$~W/kg.

\Cref{tab:margin} compares this ceiling against the ICNIRP $\psSAR$
basic restrictions.

\begin{table}[h]
\centering
\caption{Compliance margin between the APD-derived bound on $\psSAR$
and the ICNIRP $\psSAR$ basic restriction, at $\APD = 10$~W/m$^2$.}
\label{tab:margin}
\begin{tabular}{lccc}
\toprule
Body region & ICNIRP limit [W/kg] & Bound [W/kg] & Margin \\
\midrule
Head, trunk & $2.0$ & $0.465$ & $4.3\times$ \\
Limbs       & $4.0$ & $0.465$ & $8.6\times$ \\
\bottomrule
\end{tabular}
\end{table}

\section{The body-surface footprint}\label{sec:footprint}

\Cref{thm:apd-dom} uses footprint $A = L^2$. This holds when the cube
z-axis aligns with the outward body normal $\nhat$, which is the natural
alignment in the planar three-layer model: the body is a stack of plane
layers and the cube sits with one face parallel to those layers.
Generalizing to cubes whose axes do not align with $\nhat$ requires
care.

\subsection{Maximum cube cross-section by a plane}

For an axis-aligned unit cube cut by a plane through its centroid, the
cross-section area depends on the plane's normal direction $\nhat$:
\begin{itemize}\setlength{\itemsep}{2pt}
  \item $\nhat$ along a face normal (e.g., $\hat{z}$): the cut is the
        unit square, area $1$.
  \item $\nhat$ along a face diagonal (e.g., $(1,1,0)/\sqrt{2}$): the
        cut is a rectangle with sides $1$ and $\sqrt{2}$, area
        $\sqrt{2} \approx 1.414$.
  \item $\nhat$ along the body diagonal $(1,1,1)/\sqrt{3}$: the cut is
        a regular hexagon with side $1/\sqrt{2}$, area
        $3\sqrt{3}/4 \approx 1.299$.
\end{itemize}
The maximum cross-section is $\sqrt{2}\,L^2$, achieved by the
face-diagonal cut~\cite{KuhnelMath}. The body-diagonal hexagon is
smaller despite its higher symmetry.

\subsection{What the IEC mass rule does to the maximum}

The $\sqrt{2}\,L^2$ maximum occurs for the centered cut, where the
cube has $50\%$ of its volume on each side of the body plane. That
placement violates the IEC $90\%$-tissue rule. Pushing the cube deeper
into tissue to recover the $90\%$ rule moves the cut from the cube
center toward one face or edge, and the intersection area shrinks.
The realistic maximum under the IEC constraint sits between $L^2$ and
$\sqrt{2}\,L^2$, approaching $L^2$ for cubes whose axes are nearly
aligned with $\nhat$.

\subsection{Generalized bound}

For a cube whose axes are not aligned with $\nhat$, the bound carries a
geometric prefactor:
\begin{equation}\label{eq:apd-dom-gen}
  \psSAR \;\le\; \frac{\sqrt{2}\,\APD}{\rhom\,L}\, .
\end{equation}
Numerically, the worst case at the ICNIRP $\APD$ limit is
$\psSAR \le 0.658$~W/kg, still below the head and trunk restriction
($2$~W/kg) by a factor of $3.0$ and the limb restriction ($4$~W/kg) by
a factor of $6.1$. The conclusion that APD compliance certifies
$\psSAR$ compliance survives the worst-case tilt.

\section{Tightness}\label{sec:tightness}

\subsection{The capture fraction}

\Cref{eq:apd-dom} is tight when all flux entering through the footprint
deposits inside the cube. For a homogeneous absorbing half-space with
amplitude attenuation $\alpha$, the SAR depth profile is
$\SAR(z) = \SAR(0)\,e^{-2\alpha|z|}$, and energy conservation pegs the
surface SAR to the absorbed flux:
\begin{equation}\label{eq:SAR0}
  \rhom \int_{-\infty}^{0} \SAR(z)\,\dd z
  = \frac{\rhom\,\SAR(0)}{2\alpha}
  = \APD\, ,
\end{equation}
giving $\SAR(0) = 2\alpha\,\APD/\rhom$. The cube average for a face-flush
placement (top at $z = 0$, bottom at $z = -L$) is
\begin{equation}\label{eq:cube-avg}
  \langle\SAR\rangle_{\mathrm{cube}}
  = \frac{1}{L}\int_{-L}^{0} \SAR(z)\,\dd z
  = \frac{\SAR(0)}{2\alpha L}\,(1 - e^{-2\alpha L})
  = \frac{\APD}{\rhom L}\,(1 - e^{-2\alpha L})\, .
\end{equation}
The factor $C(\alpha L) = 1 - e^{-2\alpha L}$ is the capture fraction:
the share of the entering flux that deposits inside the cube depth.

\subsection{Numerical capture fractions}

\Cref{tab:capture} lists the capture fraction on skin across frequency,
using IT'IS Cole--Cole permittivity~\cite{ITIS}. Above $5$~GHz the
bound is exact to within a few percent. Below $1$~GHz the bound is
loose by $\sim 50\%$.

\begin{table}[h]
\centering
\caption{Field amplitude decay rate $\alpha$ on skin and the
corresponding capture fraction $C(\alpha L) = 1 - e^{-2\alpha L}$ for
the $L = 21.5$~mm cube.}\label{tab:capture}
\begin{tabular}{rccc}
\toprule
$f$ [GHz] & $1/(2\alpha)$ [mm] & $2\alpha L$ & $C$ \\
\midrule
0.9  & 19.9 & 1.1  & 0.67 \\
2.45 & 11.3 & 1.9  & 0.85 \\
5    &  5.3 & 4.1  & 0.98 \\
10   &  1.9 & 11.3 & 1.00 \\
28   & 0.46 & 46.7 & 1.00 \\
60   & 0.24 & 89.9 & 1.00 \\
\bottomrule
\end{tabular}
\end{table}

The capture fraction degrades on fat below $\sim 3$~GHz, where the field
penetration depth exceeds the cube side and a substantial portion of
the absorbed power deposits in muscle below the cube. In the
fat-resonance regime the layered $\Tlay$ replaces $\Tzero$ in
$\APD$, but the bound itself remains valid because it depends only on
the energy entering the cube, not on the depth profile.

\subsection{Numerical validation on Thelonious}

A 3-layer skin/fat/muscle Chew-recursion model with Monte Carlo cube
sampling on the Thelonious mesh~\cite{ThelonioPhantom} returns a median
ratio of bound-to-numerical of $1.30$ and a maximum of $1.46$ across
$1$--$28$~GHz, four polarization configurations, and five incidence
directions per surface point. The median is set by the capture fraction
of \cref{tab:capture} weighted across the swept frequencies; the
maximum is set by the worst-case frequency-incidence combination.
The validation confirms \cref{eq:apd-dom} in the realistic anatomy
regime.

\section{Implications}\label{sec:implications}

\subsection{Compliance certification through $\APD$ alone}

\Cref{thm:apd-dom} reduces $\psSAR$ compliance to $\APD$ compliance
with a fixed margin. Anywhere the closed-form chain delivers
$\APD \le 10$~W/m$^2$, the cube quantity is automatically below
$0.66$~W/kg (worst case) and well below the ICNIRP basic restrictions
on $\psSAR$. No volumetric cube search is required for compliance
certification.

This collapses the operational compliance check to a single surface
quantity. Across $1$--$100$~GHz the regulator evaluates one closed-form
expression per body-surface point. The cube search remains useful for
verification or for the small set of cases where the actual
$\psSAR$ value is required (rather than a compliance certificate), but
plays no role in the certification path.

\subsection{Connection to the closed-form chain}

The closed-form chain of the main paper~\cite{TAP_paper} delivers
$\APD$ at any body surface point as
\begin{equation}\label{eq:APD-chain}
  \APD(\mathbf{r}) = \Sinc(\mathbf{r}) \cdot \cos\theta_i(\mathbf{r}) \cdot
  \Tlay\bigl(f, \theta_i(\mathbf{r}), \mathrm{pol}\bigr)\, ,
\end{equation}
with $\Tlay$ the three-layer tissue transmission, $\theta_i$ the local
incidence angle, and $\cos\theta_i = -\nhat \cdot \khat$. The $\psSAR$
bound \cref{eq:apd-dom} composes directly:
\begin{equation}\label{eq:psSAR-chain}
  \psSAR(\mathbf{r}) \;\le\; \frac{\Sinc(\mathbf{r})\,\cos\theta_i\,
  \Tlay(f, \theta_i, \mathrm{pol})}{\rhom\,L}\, .
\end{equation}
The same surface integral that delivers whole-body absorbed power via
the generalized Cauchy identity, and $\APD$ pointwise via the geometric
local law, also delivers a $\psSAR$ compliance certificate via
\cref{eq:psSAR-chain}. Three regulatory outputs from one identity.

\subsection{Regulatory framing}

ICNIRP retains $\psSAR$ as the basic restriction below $6$~GHz and
switches to $\APD$ above $6$~GHz~\cite{ICNIRP2020}. \Cref{thm:apd-dom}
shows that an $\APD$-based certification path is sufficient across the
whole band: the basic-restriction switch at $6$~GHz is a regulatory
convention that reflects historical metric choice and the relative
maturity of measurement techniques, not a physical necessity. Both
restrictions can be certified through the same surface quantity.

This does not say the $\psSAR$ basic restriction should be removed.
$\psSAR$ encodes deeper-tissue heating through the cube depth that
$\APD$ does not. The corresponding limit values were calibrated against
local temperature-rise models that include this depth contribution. The
result here is narrower: \emph{at the limits ICNIRP currently
specifies}, satisfying $\APD \le 10$~W/m$^2$ guarantees
$\psSAR \le 0.66$~W/kg, which is well within the
$2$~W/kg head and trunk restriction. If ICNIRP tightens
$\APD$ or relaxes $\psSAR$, the dominance margin will shift; the
inequality \cref{eq:apd-dom} is unchanged.

\subsection{What the bound does not say}

Three caveats:

\begin{enumerate}\setlength{\itemsep}{4pt}
  \item \textbf{Whole-body SAR is a separate restriction.} The
        ICNIRP $\mathrm{SAR}_{\mathrm{wb}} \le 0.08$~W/kg basic
        restriction is governed by the generalized Cauchy formula on
        the body surface integral, not by the cube bound. The cube
        bound speaks only to the local restriction.

  \item \textbf{The cube footprint is not the ICNIRP $\APD$ averaging
        area.} ICNIRP averages $\APD$ over $4$~cm$^2$ below $30$~GHz
        and $1$~cm$^2$ above. The cube footprint is
        $L^2 = 4.62$~cm$^2$ for the head and trunk cube, slightly
        larger than the $\APD$ averaging area. The peak $\APD$ over
        $4$~cm$^2$ is at most $\sim 16\%$ above the average over the
        cube footprint when the field is sharp, which preserves the
        margin in \cref{tab:margin} to within rounding.

  \item \textbf{Planar three-layer model is the approximation.} On
        sharp body features (curvature radius $\lesssim L$), the
        cube intercepts a non-planar body patch and the footprint
        analysis of \cref{sec:footprint} requires correction. The
        $\sqrt{2}$ slack in \cref{eq:apd-dom-gen} absorbs typical
        curvature effects, but pathological geometries (e.g., cubes
        on body extremities thinner than $L$) are outside the bound's
        scope.
\end{enumerate}

\section{Summary}\label{sec:summary}

The inequality
\begin{equation*}
  \psSAR \;\le\; \frac{\APD}{\rhom L}
\end{equation*}
follows from energy conservation on the IEC/IEEE 62704-1 cube. The
ICNIRP $\APD$ basic restriction ($10$~W/m$^2$) implies
$\psSAR \le 0.47$~W/kg, below the corresponding $\psSAR$ basic
restrictions ($2$ and $4$~W/kg) by factors of $4.3$ and $8.6$. The
bound is tight to within a capture-fraction factor that approaches
unity above $5$~GHz and reaches $0.67$ at $0.9$~GHz on skin. Numerical
validation on the Thelonious phantom confirms a median bound ratio of
$1.30$ and a maximum of $1.46$.

The operational consequence: the closed-form chain that delivers
$\APD$ via the layered transmission $\Tlay$ certifies $\psSAR$
compliance through the same surface integral. ICNIRP's two local
basic restrictions reduce to one closed-form check across $1$--$100$~GHz.

\bibliographystyle{plain}
\begin{thebibliography}{9}
\bibitem{ICNIRP2020}
ICNIRP, ``Guidelines for limiting exposure to electromagnetic fields
($100$~kHz to $300$~GHz),'' \emph{Health Physics}, vol.~118, no.~5,
pp.~483--524, 2020.

\bibitem{62704-1}
IEC/IEEE 62704-1, ``Determining the peak spatial-average specific
absorption rate (SAR) in the human body from wireless communications
devices, $30$~MHz to $6$~GHz,'' IEC/IEEE International Standard, 2017.

\bibitem{ITIS}
P.~A. Hasgall, F. Di Gennaro, C. Baumgartner, E. Neufeld, M.~C. Gosselin,
D. Payne, A. Klingenb\"ock, and N. Kuster, ``IT'IS Database for thermal
and electromagnetic parameters of biological tissues, version 4.1,''
DOI: 10.13099/VIP21000-04-1, 2022.

\bibitem{KuhnelMath}
W. K\"uhnel, \emph{Differential Geometry: Curves -- Surfaces --
Manifolds}, AMS, 2nd ed., 2006. (Cube cross-section maxima: classical
result, also worked out in undergraduate solid geometry references.)

\bibitem{ThelonioPhantom}
M.-C. Gosselin et al., ``Development of a new generation of high-resolution
anatomical models for medical device evaluation: the Virtual Population
3.0,'' \emph{Phys. Med. Biol.}, vol.~59, no.~18, pp.~5287--5303, 2014.
(Thelonious is the 6-year-old male of the Virtual Family.)

\bibitem{TAP_paper}
R. Wydaeghe, ``Closed-form absorbed-power dosimetry from $1$ to
$100$~GHz,'' draft, 2026.

\end{thebibliography}

\end{document}
